\setauthor{Halilovic Ema}
Bei der Auswahl der Technologien für den serverseitigen Bereich wurden folgenden Kriterien berücksichtigt:
Die Auswahl der Technologien für den serverseitigen Bereich berücksichtigt die Erfüllung folgender Kriterien:
\begin{compactitem}
    \item Sie sollen auf dem neuesten Stand der Technik sein
    \item Ausreichende Dokumentation mit ständiger Weiterentwicklung verfügen
    \item Eine weite Auswahl an \glspl{fw} und Funktionalitäten unterstützen in den Bereichen:
    \begin{compactitem}
        \item REST
        \item Filemanagement
        \item Datenbank
    \end{compactitem}
    % \item Und schon erste Erfahrungen in der dazugehörigen Programmiersprache
\end{compactitem} 
% TODO listin gwas die Frameworks machen können in welchem bereich

Neben der Erfüllung der oben genannten Kriterien wurden bei der Entscheidung zur Verwendung der Technologien auch bereits vorhandene Praxiserfahrungen berücksichtigt. 
Eine Übersicht sowie nähere Beschreibung der für die vorliegende Arbeit verwendeten Technologien sind in den folgenden Unterkapiteln gegeben.
Allerdings wurden die endgültigen Technologien aufgrund schon vorhandener Praxiserfahrung getroffen. 
% Dazu finden sich mehrere \glspl{fw}, welche in Fragen kommen.

\section{PostgreSQL}
\begin{wrapfigure}{r}{0.3\textwidth}
  \begin{center}
      \includegraphics[width=0.2\textwidth]{pics/logos/postgres.png}
      \caption{PostgreSQL Logo}
  \end{center}
\end{wrapfigure}
\emph{PostgreSQL} ist ein Open Source System für relationale Datenbanken. 
Es wurde 1986 erstmals an der University of Califorinia at Berkeley entwickelt, was eine Entwicklung seit 35 Jahren bedeutet.  
\emph{PostgreSQL} läuft auf den gängigsten Betriebssystemen im Terminal.
Durch die Entwicklung des Systems seit rund 35 Jahren, entstanden ebenso mehrere User-Interfaces, wie beispielsweise PgAdmin. 
\cite{PostgreSQLAbout}

PgAdmin ist das am weitesten verbreitete Administrationstool für PostgreSQL. 
\cite{PgAdminAbout}
Es erleichtert die Ausführung von \gls{sql}-Befehlen, sowie den Wechsel verschiedener Tabellen.

Eine gute Dokumentation und schon vorhandene Praxiserfahrung haben diese Datenbank hervorgehoben und zur endgültigen Entscheidung geführt. 
In dieser Arbeit wurde PgAdmin 4 verwendet, um Daten mittels Port-forwarding auf der Cloud zu überprüfen. 
Durch die Visualisierung der Tabellen war diese Aufgabe einfacher als mit dem Terminal allein. 

\section{Quarkus}
\begin{wrapfigure}{r}{0.3\textwidth}
    \begin{center}
      \includegraphics[width=0.2\textwidth]{pics/logos/quarkus.png}
     \caption{Quarkus Logo}
    \end{center}
  \end{wrapfigure}
Das Open Source \gls{fw} Quarkus wird für die Entwicklung von cloud-native Projekten in Java verwendet. 
Es zeichnet sich unter anderem durch die kurzen Startzeiten, sowie den geringen Arbeitsspeicherverbrauch aus. 
Ebenso sind Quarkus-Projekte dazu optimiert, in Containerumgebungen, wie beispielsweise Kubernetes, ausgeführt zu werden.
\cite{QuarkusHomepage, QuarkusAbout}

Ein neues Projekt wird entweder durch das Quarkus-\gls{cli} erstellt oder nach individueller Konfiguration auf \href{https://code.quarkus.io}{code.quarkus.io} heruntergeladen. 
Standardmäßig sind in dem Projektordner \hyperref[ch::MavenTool]{Maven}-Verzeichnis, sowie Quarkus-Dateien enthalten, welche Projekteinstellungen enthalten. 
Eine dieser Konfigurationsdateien wäre zum Beispiel die \emph{application.properties}-Datei. 
\cite{QuarkusAbout, QuarkusFirstApplication}

Zusätzlich verfügt Quarkus über eine Vielzahl von Extensions, welche eine Vielzahl der gängigsten Technologien unterstützen.
Sie werden durch das \gls{cli} oder manuell hinzugefügt. 
\cite{QuarkusAbout, QuarkusFirstApplication}

Mit Quarkus wurde der Backend-Teil der Arbeit erstellt. 
Die Möglichkeit das Projekt mit zusätzlichen Erweiterungen auszustatten war 

\subsection{Maven}
\label{ch::MavenTool}
\begin{wrapfigure}{r}{0.3\textwidth}
    \begin{center}
        \includegraphics[width=0.1\textwidth]{pics/logos/asf.jpg}
        \caption{Apache feather Logo}
    \end{center}
\end{wrapfigure}

Maven ist ein Open Source Build-Tool, welches von der Apache Software Foundation lizensiert ist.
Es wurde in Java geschrieben und wird meist in Java-Projekten eingesetzt. 
Durch zentrale Nutzung des \gls{pom} werden projektrelevante Informationen gesichert, wodurch der Kompilierungsprozess vereinfacht wird.
\cite{MavenAbout, MavenAbout2gls}

Maven sorgt dafür, dass unterschiedliche Geräte eine Applikation ausführen können, ohne manuelle Konfiguration. 
Die einzige Voraussetzung  des Zielgerätes ist, dass Maven installiert und eingerichtet ist. 
Ebenso können \emph{Maven-Ordner} gedownloadet werden. 
Diese ermöglichen die Nutzung der Maven Befehle ohne Installation des Tools am Gerät.
\cite{MavenAbout, MavenAbout2gls}

Quarkus-Projekte verwenden Mavens \emph{pom.xml}-Datei, um unter anderem die verwendete Java Version oder die verwendeten Extensions abzuspeichern.
Zusätzlich wird ein einheitliches System für Projektkonfigurationen geboten. 
Dadurch müssen, wie oben erwähnt, Einstellungen bei Gerätewechsel nicht mehr manuell geändert werden. 
\cite{MavenAbout}

\subsection{\gls*{jdbc} Treiber}

In Java wird für Datenbankverbindungen die \gls{jdbc} \gls{api} verwendet. 
Genauso benötigen Quarkus-Projekte einen datenbankspezifischen \gls{jdbc}-Treiber für einen Zugriff auf die Datenbank. 
In diesem Projekt wurde PostgreSQL als Datenbank gewählt, weshalb die Extension \emph{\gls{jdbc} Driver - PostgreSQL} fundamental für die Entwicklung war. 
\cite{datasourcesAbout}.

Um die Extension zu verwenden und eine Datenbankverbindung aufzubauen, müssen in den \emph{application.properties} zusätzliche Konfigurationen eingefügt werden. 
Wichtig sind Informationen, wie die Art der Datenbank, der Pfad, um diese zu erreichen, und die Login-Daten eines berechtigten Users.
Dazu wird der \gls{boilercode} aus Codeblock \ref{lst:quarkusDatasource} verwendet, wobei die Klammern durch tatsächliche Werte ersetzt werden. 
\cite{datasourcesAbout}

\begin{lstlisting}[caption=Beispielkonfigurationen,label=lst:quarkusDatasource]
  quarkus.datasource.db-kind=postgresql 
  quarkus.datasource.username=<meinUser>
  quarkus.datasource.password=<meinPassword>
  quarkus.datasource.jdbc.url=jdbc:postgresql://<URL>:<Port>/<meinName>
\end{lstlisting}

\subsection{Hibernate ORM mit Panache} 
Da Java eine objekt-orientiert Programmiersprache ist, und PostgreSQL eine relationale Datenbank, können standardmäßig keine Entitäten aus Java in die Datenbank übertragen werden. 
Diese Aufgabe übernimmt ein \gls{orm}, der den Java-Code für PostgreSQL umwandelt. 
Dadurch können Java-Klassen als Objekte persistiert und gelöscht werden, ohne dass komplexer Code konstruiert werden muss. 
Hibernate ORM ist die dafür zuständige Extension für Quarkus. 
\cite{ORMAbout}

Panache bietet zusätzliche Klassen mit Funktionen, von denen abgeleitet werden kann. 
Diese erleichtern das Arbeiten mit angelegten Entitäten, besonders bei \gls{crud}-Operationen.
Zum Beispiel ist für die Persistenz eines neuen Objekts lediglich ein Aufruf der \emph{.persist()}-Methode notwendig.
\cite{HibernateORMwithPanache}

\subsection{REST-Easy}
\label{ch::MavenTool}
\begin{wrapfigure}{r}{0.3\textwidth}
    \begin{center}
        \includegraphics[width=0.2\textwidth]{pics/logos/resteasy.png}
        \caption{REST Easy Logo}
    \end{center}
\end{wrapfigure}
\gls{rest}-Easy ist eine Quarkuserweiterung, die es ermöglicht, mit RESTful Web Servies zu arbeiten. 
Sie implementiert die Jakarta RESTful Web Services der Eclipse Foundation, sowie die MicroProfile \gls{rest} Client Spezifikation \gls{api}.
\cite{ResteasyAbout}

%\gls{rest} hat sich als Standard für Mikroservice-Anwendungen festgelegt.\cite{RESTAbout}
Dies bedeutet, dass durch diese Extension im Projekt \gls{api}s erstellt werden können. 
Ein Endpoint wird durch Definition zweier Annotationen \emph{@Path} und \emph{@<beliebige HTTP-Methode>} erzeugt.
\cite{ResteasyExtensionAbout}

Folgende \gls{http}-Methoden werden in dieser Arbeit verwendet:
\begin{compactitem}
    \item \textbf{GET}: GET-Requests liefern Daten \cite{httpAbout}
    \item \textbf{POST}: POST-Requests übermitteln Daten wie Entitäten \cite{httpAbout}
    \item \textbf{DELETE}: DELETE-Requests löschen Daten \cite{httpAbout}
\end{compactitem}

Die Beschreibungen der oben angeführten Methoden sind Normen in der Entwicklung mit \gls{http}. 
Eventuell hätte noch \emph{PUT} verwendet werden können, jedoch wird diese Methode für einzelne Änderungen einer Entität verwendet.
Dadurch hatte sie keinen spezifischen Nutzen in dieser Arbeit. 
\cite{httpAbout}

\subsection{Swagger-ui}
\begin{wrapfigure}{r}{0.3\textwidth}
    \begin{center}
        \includegraphics[width=0.2\textwidth]{pics/logos/swagger.png}
        \caption{Swagger Logo}
    \end{center}
\end{wrapfigure}
Swagger UI ist ein Tool zum Testen von \gls{rest}-\gls{api}s. 
Es ist Open Source und liefert mehrere Funktionalitäten, wobei in dieser Arbeit hauptsächlich die Visualisierung der Schnittstellen verwendet wurde. 
Dieser Teil bietet, je nach definierten Parametern, vorgefertigte Requests, welche noch bearbeitet werden können. 
Ebenso zeigen sie die benötigte \gls{http}-Methode und ermöglichen es, den Endpoint mit Mausklick zu testen. 
Dies erspart sehr viel Arbeit während der \gls{api}-Entwicklung.
\cite{SwaggeruiAbout}

\subsection{JUnit 5}
\begin{wrapfigure}{r}{0.3\textwidth}
    \begin{center}
        \includegraphics[width=0.2\textwidth]{pics/logos/junit5-logo.png}
        \caption{JUnit 5 Logo}
    \end{center}
\end{wrapfigure}
JUnit steht für Java Unit und ist ein Open Source \gls{fw} zum Testen. 
Es hat sich als der Standard dieser Programmiersprache bewiesen und ist spezialisiert auf die Überprüfung von Methoden und Klassen.
Durch die Popularität besteht eine ausführliche Dokumentation und weites Funktionsspektrum.
Ebenso inspirierte dieses Konzept die \gls{fw}s anderer Programmiersprachen im Testing-Bereich.

JUnit ist in vielen der gängigsten \gls{ide}s inkludiert. 
Zum Anlegen einer Testklasse wird gewöhnlich der Name der zu testenden Klasse verwendet mit \emph{.test} als Suffix. 
Die benötigten Repositories werden simuliert, sodass Methoden aus diesem wie gewohnt verwendet werden können.
\cite{JUnitAbout}


\section{IntelliJ IDEA}
\begin{wrapfigure}{r}{0.3\textwidth}
  \begin{center}
      \includegraphics[width=0.2\textwidth]{pics/logos/intellij.png}
      \caption{IntelliJ Logo}
  \end{center}
\end{wrapfigure}
IntelliJ IDEA ist eine \gls{ide}, welche von JetBrains entwickelt wurde. 
Diese ist ausgelegt für Java- und Kotlin-Projekte und assistiert bei verschiedensten \gls{jvm}-\glspl{fw}. 
Durch eingebaute Features erleichtert diese Entwicklungsumgebung das Programmieren für Nutzende. 
Plug-Ins ermöglichen es, Datenbankverbindungen und weiteres in der \gls{ide} zu konfigurieren, sodass dem Entwickler oder der Entwicklerin eine Übersicht von benötigten Informationen gegeben werden kann.
\cite{IntelliJIDEA}

IntelliJ ist nicht Open Source und benötigt ein kostenpflichtiges Abonnement. 
Als Schüler und Schülerin kann nach der Angabe von Schulinformationen eine gratis Lizenz für 6 Monate alle JetBrains Produkte beantragt werden. 
Wäre der Benefit nicht vorhanden, hätten wir die \gls{ide} nicht gewählt, sondern Visual Studio Code \ref{ch::vsc} für den ganzen Entwicklungsprozess verwendet. 
\cite{IntelliJIDEAPricing}

\section{Visual Studio Code}
\label{ch::vsc}
\begin{wrapfigure}{r}{0.3\textwidth}
  \begin{center}
      \includegraphics[width=0.2\textwidth]{pics/logos/vscode.jpg}
      \caption{VS Code Logo}
  \end{center}
\end{wrapfigure}
Visual Studio Code ist ein Open Source Code Editor für verschiedene Zwecke. 
Das Programm ist auch unter dem Namen \emph{VS Code} bekannt und wurde von Microsoft lizensiert. 
Der offen verfügbare Quellcode führt dazu, dass die Community weiterhin Erweiterungen entwickelt.
\cite{vscodeAbout}

Der Editor läuft auf den Betriebssystemen Windows, macOS und Linux. 
Bei der Installation wird anfangs nur ein Language Support für JavaScript, TypeScript und Node.js mitgeliefert. 
Durch die vielen Erweiterungen kann sich VS Code individuell angepasst werden. 
Ebenso kann es als \gls{ide} genutzt werden, wenn eine Extension für die gewünschte Sprache gefunden werden kann. 
\cite{vscodeAboutGs}

Der Editor wurde für die letzten Schritte der Arbeit genutzt. 
Dazu zählen die Konfiguration der benötigten Cloud-Dateien, sowie das von Frondend und Backend. 

\section{LeoCloud}
Die LeoCloud ist ein Projekt der HTL Leonding mit Aberger Christian als Projektleiter. 
Sie ermöglicht es mittels Kubernetes, eine beliebige Anwendung auf einer Cloud laufen zu lassen. 
Um diese Cloud nutzen zu können, ist eine E-Mail-Adresse mit der HTL-Leonding-Schulsignatur notwendig. 
Da dies ein schulinternes Projekt ist, war eine Kommunikation mit den Entwicklern gegeben, das heißt, dass Fragen sofort beantwortet werden konnten. 
\cite{LeoCloudAbout}

Die Cloud wurde für die Anforderungen genutzt, dass Projekt für alle Geräte zugänglich zu machen. 
Durch ein Deployment auf die LeoCloud war die Anwendung über den Browser erreichbar und somit ein endgültiges Produkt. 

\subsection{Kubernetes}
\begin{wrapfigure}{r}{0.3\textwidth}
  \begin{center}
    \includegraphics[width=0.2\textwidth]{pics/logos/k8s.png}
   \caption{Kubernetes Logo}
  \end{center}
\end{wrapfigure}
Kubernetes ist eine Anwendung, die für das Containern von Applikationen entwickelt wurde. 
Sie ist Open Source und weit verbreitet. 
Ebenso basiert sie auf den Produktionworkloads von Google.
Alternativ ist Kubernetes unter dem Namen \emph{k8s} bekannt.
\cite{k8sAbout}

Wie vorhin erwähnt, nutzt die LeoCloud Kubernetes, um Anwendungen auf die Cloud zu laden. 
Besonders wird dabei das \gls{cli} verwendet. 

Da das Arbeiten am Terminal unübersichtlich werden kann, wurde ebenso die Extension \emph{Kubernetes von Microsoft} in Visual Studio Code genutzt.
\cite{k8sAboutExtension}

\section{GitHub}
\begin{wrapfigure}{r}{0.3\textwidth}
  \begin{center}
      \includegraphics[width=0.2\textwidth]{pics/logos/github.png}
      \caption{GitHub Logo}
  \end{center}
\end{wrapfigure}
GitHub ist eines der am weitest verbreiteten Entwicklertools auf der Welt. 
Es ist Open Source verfügbar und ermöglicht mit einer Webseite // TODO 
Zusätzlich existieren bereits diverse Apps und Extensions, die die Entwicklung mit GitHub auf anderen Plattformen verbreiten.
\cite{githubAbout}

Durch GitHub Repositories wird das Konzept der \gls{ssot} ermöglicht. 
Bei der Erstellung von Repositories werden mehrere Funktionen geboten, wie Pull Requests. 
Wenn ein Teammitglied eine Änderung des Codes veröffentlicht, überschreiben diese nicht den aktuellen Stand, sondern erzeugen Pull Requests.
Nach einer Überprüfung dank der integrierter Codevorschau können die Requests geschlossen und die neuen Zeilen mit schon vorhandenem Code zusammengeführt werden. 
Alle Änderungen bleiben in Changelogs gespeichert, sodass immer auf alte Versionen zugegriffen werden kann.
\cite{githubAbout}

\section{Notion}
\begin{wrapfigure}{r}{0.3\textwidth}
  \begin{center}
    \includegraphics[width=0.2\textwidth]{pics/logos/notion.png}
   \caption{Notion Logo}
  \end{center}
\end{wrapfigure}

Notion ist eine Art digitales Notizbuch, womit so genannte \emph{Workspaces} erstellt werden können. 
Es ist im Browser, als App in Windows-, MacOS- oder auf iOS- und Android-Geräten verfügbar. 
Workspaces aufgebaut durch mehrere Seiten, welche durch Befehle erstellt werden. 
Diese werden auch direkt miteinander verlinkt.
Seiten lassen sich mit beliebigem Inhalt füllen und durch einfache \emph{/-Befehle} ist es möglich z. B. ein Kanban Board erstellen. 
\cite{NotionAbout}

Notion ist gut für Neueinsteiger, da alle Funktionalitäten gut beschrieben sind und es benutzerfreundlich gestaltet ist.
Workspaces lassen sich mit anderen teilen, wodurch jeder User auf eine \gls{ssot} Version des Workspaces zugreifen kann.
Für die Bearbeitung in echt-zeit wird Internet benötigt, falls dieses jedoch ausfällt, werden die Änderungen gespeichert und synchronisiert, sobald wieder eine Netzwerkverbindung aufgebaut wurde.
\cite{NotionAbout}

Das Feature von geteilten Workspaces wird in dieser Arbeit als Erleichterung der Kommunikation zwischen Teammitgliedern genutzt, um gemeinsame Notizen und Termine festzuhalten. 
Durch die Verwendung dieses Tools wurde sichergestellt, dass jedes Mitglied eine \gls{ssot}-Version der geteilten Informationen hatte.
Für die selbstständige Arbeit wurde das Tool ebenso genutzt, um wichtige Webseiten zu speichern, sowie Zwischennotizen zu erstellen. 

\section{PlantUML}
\begin{wrapfigure}{r}{0.3\textwidth}
  \begin{center}
      \includegraphics[width=0.2\textwidth]{pics/logos/plantuml.png}
      \caption{PlantUML Logo}
  \end{center}
\end{wrapfigure}
PlantUML ist ein Tool, welches die Erstellung von \gls{uml}-Diagrammen vereinfacht. 
Das Ziel von \gls{uml} ist es, Tools zu liefern, um Systeme zu visualisieren.
\cite{UMLPaper}
Dafür gibt es verschiede Arten von \gls{uml}-Diagrammen, wie beispielsweise Objektdiagramme oder Sequenzdiagramme.
\cite{PlantUML}

Für diese Arbeit wird die gratis Extension \emph{PlantUML von jebbs} verwendet. 
Sie bietet eine Preview des Diagramms, sowie Exportmöglichkeiten. 
\cite{PlantumlExtension}
Die Visualisierung unserer Datenstruktur geschieht durch die Erstellung eines \gls{erd}s mit den PlantUML Klassendiagrammen. 
Die Beziehungen zwischen den Entitäten werden hier mittels der Krähenfußnotation dargestellt.

\section{JSON Web Token}
\begin{wrapfigure}{r}{0.3\textwidth}
  \begin{center}
      \includegraphics[width=0.2\textwidth]{pics/logos/jwt.png}
      \caption{JWT Logo}
  \end{center}
\end{wrapfigure}
JSON Web Token, oder \gls{jwt}, ist ein, durch RFC 7519 gekennzeichneter, Standard für die sichere Datenübertragung. 
Die Token übermitteln Daten als digital signierte \gls{json}-Objekte.
Die Signatur erfolgt durch eine Verschlüsselung mittels Secret oder Schlüsselpaar. 
Secrets werden mit einem HMAC-Algorithmus erstellt, wobei Schlüsselpaare RSA oder ECDSA verwenden
\cite{JWTAbout}

In Tokens kann eine freie Anzahl von wählbare Daten, wie selbst definierte Adjektive, gespeichert werden. 
Typischerweise werden Tokenart und Verschlüsselungsalgorithmus im Header des Tokens gespeichert und die individuell definierten in der \emph{Payload}.
Der Token wird bei einem erfolgreichem Login Request erhalten und danach bei \gls{http}-Requests im Authorization Header mitgeschickt. 
Damit der Token nicht unendlich lang genutzt werden kann, ist es wichtig, eine \gls{ttl} festzulegen. 
Bei Interesse an den gespeicherten Informationen kann jeder JWT auf \href{https://jwt.io}{https://jwt.io} enkodiert werden.
\cite{JWTAbout}

\section{Homebrew}
\begin{wrapfigure}{r}{0.3\textwidth}
  \begin{center}
      \includegraphics[width=0.1\textwidth]{pics/logos/brew.png}
      \caption{Homebrew Logo}
  \end{center}
\end{wrapfigure}
Homebrew ist ein Paketmanager für Linux und macOS. 
Er wurde in Ruby von Max Howell programmiert und vereinfacht Installation und Aktualisierung von Tools und Dependencies. 
Dazu müssen sie im Homebrew-Repository enthalten sein. 
Alle Installationen werden im selben Ordner gespeichert, wodurch bei Bedarf des Pfades dieser sofort gefunden wird. 
Um Homebrew zu verwenden, wird im Terminal Stichwort \emph{brew} genutzt. 

Die meisten der oben angeführten Technologien sind im Homebrew-Repository enthalten. 
Durch die Entwicklung auf einem macOS-System, war dieser Paketmanager ein sehr hilfreiches Tool. 
\cite{brewAbout}

%\subsection{Deeper}
%Nicht mehr im Inhaltsverzeichnis.

%\subsubsection{Deepest}
%Vermeide mich.